\section{推理策略：从 Few-shot 到 Zero-shot}

\subsection{Least-to-Most Prompting：分解复杂任务}

受 P\'{o}lya 经典著作《How to Solve It》中分解策略的启发，Denny Zhou 团队提出了 Least-to-Most Prompting：

\begin{itemize}
    \item \textbf{核心思想}：先将复杂问题分解为一系列子问题，然后从最简单的子问题开始逐步求解，每个子问题的解作为后续子问题的上下文。
    \item \textbf{SCAN 任务}：组合泛化（compositional generalization）基准，仅用 1\% 的示例即达到 99.7\% 准确率。
    \item \textbf{Text-to-Code 任务}：使用 Dynamic Least-to-Most Prompting，仅用 1\% 的数据即大幅超越使用全部训练数据的专用架构。
\end{itemize}

\begin{knowledgebox}{组合泛化（Compositional Generalization）}
指测试样本比训练/提示样本更复杂的场景。例如，训练时只见过短代码片段，测试时需要生成更长的代码。Least-to-Most Prompting 通过分解策略天然地处理了这一问题。
\end{knowledgebox}

\subsection{Zero-shot CoT：``Let's think step by step''}

\begin{itemize}
    \item 无需任何示例，仅在问题后加一句 ``Let's think step by step'' 即可触发 LLM 的逐步推理。
    \item 效果通常不如 few-shot CoT，但胜在零成本。
\end{itemize}

\subsection{LLMs as Analogical Reasoners}

受 P\'{o}lya 的类比推理思想启发：

\begin{itemize}
    \item \textbf{方法}：面对新问题时，让 LLM 自行生成相关问题及其解法，再利用这些自生成的示例来解决原问题。
    \item \textbf{优势}：比 ``Let's think step by step'' 显著更好，甚至超越手工设计的 few-shot 示例——因为模型会为每个问题生成定制化的相关示例。
    \item \textbf{与检索增强的关系}：可以进一步扩展为从网络检索相关问题和知识，实现 scaling。
\end{itemize}

\subsection{Chain-of-Thought Reasoning without Prompting}

一个令人惊讶的发现：\textbf{无需任何提示}，仅通过特殊解码策略即可触发推理：

\begin{itemize}
    \item 在解码第一步，不取概率最高的 token，而是探索 top-$k$ 个候选 token。
    \item 对每个候选 token 继续贪心解码，生成完整回答。
    \item 选择包含推理步骤且最终答案置信度最高的路径。
\end{itemize}

\begin{importantbox}{关键观察}
LLM 内部已经"知道"如何推理。当推理路径（如 ``Nicolas Cage was born in 1964, and 1964 is an even year''）出现时，最终答案的概率从低值跃升至 98\%。推理步骤在概率层面显著提升了模型对最终答案的置信度。
\end{importantbox}

\subsection{本章小结}

从 few-shot CoT 到 zero-shot CoT，再到无需提示的特殊解码策略，LLM 推理能力的触发方式越来越灵活。核心规律是：推理步骤本身而非触发方式才是关键。类比推理进一步表明，LLM 可以自主生成相关示例来辅助推理。

